\chapter{Gearbox}
\section{Introduction}
    Gear-driven components: turboprop engine propellers, helicopter main and tail rotors, and various engine 
    auxiliaries coordinate the speeds of power-consuming shafts with the engine rotor. Because propellers and 
    most auxiliaries run at lower speeds, a gearbox (reduction drive) is required to reduce the engine's peripheral 
    rotation. While basic gearbox design principles are taught in “Mechanics,” aircraft gearboxes possess 
    distinctive characteristics: they demand higher cooling intensity due to increased oiling; operate in 
    multi-mode, shortening their service life; face challenges maintaining precise steady modes, especially 
    for turboprop (TPE) and turboshaft (TShE) engines due to torque-metering unit difficulties; suffer 
    installation complications at the engine front, affecting velocity triangles and stressing the compressor 
    face; and must be as light as possible.

\section{Gearbox Limitations}
    The most crucial parameters of the gearbox are mass and diameter. Mass of modern TPE gearboxes makes up 
    25--30\,\% of the overall engine mass. The gearbox of a helicopter is usually two or three times heavier 
    than the TShE itself. The greater the transmitted power and the gear ratio, the heavier the gearbox is. 
    The specific mass of the TPE gearbox (ratio of the gearbox mass to the power transmitted) ranges 
    within \(0.05\text{-}0.1\,\mathrm{kg}\;\mathrm{kW}^{-1}\).

    Diametrical size must be minimal. This is of great importance for the gearbox inside the engine. 
    If the gearbox diameter exceeds the diameter of the compressor, it is very hard to guide air to axial 
    compressor inlet with minimal hydraulic losses.

    Thus, the task for the engineers is to design a light, reliable, and small gearbox able to transfer power 
    at high efficiency. The problem becomes complicated if the gearbox is expected to transmit great power 
    values or when its gear ratio is high.

    The gearboxes driving a single lift rotor or double lift rotor are designed using simple, planetary and 
    differential gears. They can be single-stage or multiple-stage depending on gear ratio.

\section{Classification of Gearbox}
    Gearboxes can be classified by different features: 
		\begin{enumerate}
			\item \textbf{Number of input drive shafts}\par
            Conventional gearboxes have one input and one output shaft. However, designers often use alternative 
            configurations: a single input with two (usually coaxial) outputs for high-power engines (e.g., NK-12); 
            a branched diagram with two outputs for helicopters with a lift and tail rotor; a gearbox with two 
            inputs and one output when two turboprop engines drive a single propeller; and a single gearbox serving 
            two engines (e.g., Mi-24's VR-24) that powers lift and tail rotors or coaxial lift rotors (Ka-25). 
            Thus, gearboxes can feature combinations such as two inputs and two outputs, as seen in the RV-3F 
            driven by GTD-3F engines.
			\item \textbf{Shaft arrangement}\par
			Turboprop engines use coaxial gearboxes. In twin-engine configurations, the gearbox uses two parallel 
            input shafts, whereas gearboxes with two intersecting shafts are used exclusively in helicopter power plants.
			\item \textbf{Number of stages}\par
			Gearboxes may comprise one, two, three, or four stages, with the four-stage configuration reserved for 
            helicopter power plants. The arrangement of the stages can vary between designs.
            \item \textbf{Kinetics of the stage}\par
            Stages of gearboxes can be grouped into three types: simple, planetary, and differential. Differential 
            stages are further divided into loop (single degree of freedom) and two-degree-of-freedom variants.
            \item \textbf{Gear type}\par
            Most gearboxes use spur gears of diverse sizes, while bevel wheel trains appear only in configurations 
            with crossing input shafts, such as the gear trains found in helicopters.
		\end{enumerate}

\section{Fundamentals for Designing the Gearbox}
    \subsection{Choose the Gear Type}
        \begin{itemize}
			\item \textbf{Spur Gears}\par
            Simple gears with straight teeth. Best for low-speed, high-torque applications.
			\item \textbf{Helical Gears}\par
            Teeth are cut at an angle, allowing for smoother and quieter operation. Used for higher speed applications.
			\item \textbf{Bevel Gears}\par
            Used for transmitting power at a 90 degree angle.
            \item \textbf{Worm Gears}\par
            A gear arrangement where a worm (a gear with a screw-like thread) meshes with a worm wheel, providing high reduction ratios and compact design.
            \item \textbf{Rack and Pinion}\par
            Used to convert rotational motion into linear motion.
		\end{itemize}

    \subsection{Determine Gear Ratio}
        The gear ratio is crucial for controlling the speed and torque at the output shaft. The formula for calculating the 
        ratio will differ depending on the type of gearbox arrangement. The formula for a worm gearbox is different for that 
        of an epicyclic gearbox.
            \begin{equation}
                \label{eq:singlestagegearratio}
                m \;=\; \frac{T_{\text{G}}}{T_{\text{P}}}
            \end{equation}

        \autoref{eq:singlestagegearratio} represents a single stage gearbox and would be ideal when designing a worm gearbox or a 
        compound gear train. A note to keep in mind is a higher gear ratio provides more torque but reduces speed, while 
        a lower ratio gives higher speed with less torque. 
            \begin{equation}
                \label{eq:triplestagegearratio}
                m \;=\;
                \frac{G_{\text{2}}}{G_{\text{1}}}
                \;\cdot\;
                \frac{G_{\text{4}}}{G_{\text{3}}}
                \;\cdot\;
                \frac{G_{\text{6}}}{G_{\text{5}}}
            \end{equation}
        
        \autoref{eq:triplestagegearratio} would be applied for a multistage compound gearbox. The three stages represent 
        the reduction between each set of shafts. Another important aspect to keep in mind is maintaining the size 
        and weight of the gearbox. For example, a total reduction ratio of 15:1 is desired to reduce the speed of 3000 rpm.
            \begin{equation}
                \label{eq:speedreduction}
                n_{\text{output}}
                \;=\;
                \frac{n_{\text{input}}}{m}
                \;=\;
                \frac{3000~\text{rpm}}{15}
                \;=\;
                200~\text{rpm}
            \end{equation}
        
        With this information, two factors are known, the input speed, and what is required for the output speed. 
        The next phase of developing the gear train is keeping each gear as small as possible while still meeting 
        the speed and strength requirements. This can be achieved for a gearbox by taking the cube root of the 
        gear ratio as shown in \autoref{eq:ratioreductionformula}. The cube root is being used because this example 
        is a three-stage gear reduction. 
            \begin{equation}
                \label{eq:ratioreductionformula}
                m_{\text{ea}}
                \;=\;
                \sqrt[3]{m}
                \;=\;
                \sqrt[3]{15}
                \;=\;
                2.47
                \;\approx\;
                2.5
            \end{equation}

        The value produced by \autoref{eq:ratioreductionformula} should be the targeted gear ratio between each of the three stages. 
        Given that it is not a whole number, the value may need to be rounded up or down depending upon the 
        output speed requirements. If a specific stage must be adjusted to a larger or smaller ratio, 
        that will be fine, but this serves as a good starting point in the development. Following this 
        procedure will allow for having the same gears in each ratio, which will mean similar shafting, 
        bearings and mounting methods.  

    \subsection{Material and Component Selection}
        The material selection for gearbox components is a critical engineering decision that must account 
        for the distinct functional requirements of each element, ranging from aluminum housings and steel 
        bearings to medium-carbon steel shafts and exposed gears. The applicability of a particular material 
        type depends on load conditions, expected service life, cost constraints, and fabrication feasibility, 
        necessitating meticulous evaluation of mechanical properties, corrosion resistance, and manufacturability 
        for each part. Common gear materials include high-strength steel, which offers excellent durability 
        for high-load applications; advanced alloys, which provide lightweight yet high-performance 
        characteristics suitable for aerospace and other demanding sectors; bronze, whose inherent 
        wear resistance makes it the material of choice in worm-gear configurations; and engineering plastics, 
        which are often deployed in lower-load or noise-sensitive environments. Each material 
        class must be carefully matched to its respective component to achieve optimal performance and 
        longevity of the overall gearbox system.

\section{Multi-stage Gear Transmission Design}
    The objective functions considered in this work are geartrain efficiency and total weight, 
    aiming to design a transmission that maximizes efficiency while minimizing mass.

    \subsection{Geartrain Weight}
        The overall volume of the system is approximated as the sum of the volumes of each gear. Rather than 
        modeling each gear tooth explicitly, the gear is represented as a disk with an outer diameter equal 
        to the pitch diameter and an inner diameter corresponding to the defined bore. In this simplified 
        approach, the volume of each gear is computed as the product of its cross-sectional area and face width. 
        The total weight of the transmission can be calculated as:
            \begin{equation}
                \label{eq:geartrainweight}
                W \;=\; \rho_{\text{steel}}\;V
                \;=\; 7.8\times10^{3}\;\text{kg/m}^{3}\;
                \sum_{i=1}^{n}\pi\,b_{\text{i}}
                \Bigg[\Big(\frac{m_{\text{i}}\,z_{\text{i}}}{2}\Big)^{2}
                \;-\; r_{\text{bore}}^{2}\Bigg]
            \end{equation}
        This assumption approximates the bulk geometry while decreasing computational effort. 

        Volumetric Approximation: In preliminary optimization, the total mass is estimated by summing the volumes 
        of each gear, modeled as disks. The formula used is: \[\text{Total Weight} \;=\;\sum \bigl(\,\text{Cross-sectional Area}\;\times\;\text{Face Width}\;\times\;\text{Material Density}\,\bigr)\]
        This approach approximates the bulk geometry to decrease computational effort during design iterations.

    \subsection{Design Constraints}
        The constraints applied in this study are divided into three categories: geometric, layout, and 
        safety constraints. Geometric constraints relate to the dimensions, form of individual gears, 
        and ensure the feasibility of each gear stage. Layout constraints address spatial and positional 
        relationships between multiple gear stages, maintaining compatibility within the overall 
        transmission architecture, and satisfying specific requirements. Safety constraints concern 
        the strength of the gears and are assessed through ISO standardized safety factors.

    \subsection{Safety Constraints}
        The root and flank safety factors are the ratios between the modified value of 
        allowable bending ($\sigma_{at}$) and contact ($\sigma_{ac}$), and the calculated bending ($\sigma_{t}$) and 
        contact ($\sigma_{c}$) stress as shown in \autoref{eq:safetyconstraints1} and \autoref{eq:safetyconstraints2}. The modification or service factors 
        are stress cycle factor ($Y_{N}$) for bending strength, which is imposed 0.7 for idler and 1 
        for simple gears, temperature factor ($Y_{\sigma}$), reliability factor ($Y_{Z}$), stress cycle factor ($Z_{N}$)
        for contact stress, and hardness ratio factor ($C_{H}$) for contact stress.
            \begin{equation}
                \label{eq:safetyconstraints1}
                S_F \;=\;
                \frac{\sigma_{\!at}\,Y_N}{\sigma_{\!t}\,Y_{\sigma}\,Y_Z}
                \;\ge\;S_{F,\min}
            \end{equation}

            \begin{equation}
                \label{eq:safetyconstraints2}
                S_H \;=\;
                \frac{\sigma_{\!ac}\,Z_N\,C_H}{\sigma_{\!c}\,Y_{\sigma}\,Y_Z}
                \;\ge\;S_{H,\min}
            \end{equation}
        
        The calculation of bending stress ($v_{t}$) and contact stress ($v_{c}$), as well as the corresponding 
        unmodified allowable stress values, follow the procedures outlined in ISO 6336. 
        These values are derived based on the applied tangential load, gear tooth geometry, 
        and material properties.